\clearpage
\section[Обзор методов решения задачи 3D Bin Packing]{\centering\MakeUppercase{Обзор методов решения задачи 3D Bin Packing}}

Задача 3D Bin Packing относится к классу NP-трудных задач упаковки, поэтому в практических системах редко используется полный перебор всех размещений. В ходе практики рассматривались методы, которые дают допустимое решение за ограниченное время и могут быть проверены единым валидатором: последовательные greedy-эвристики, послойные схемы MaxRects, метаэвристика Large Neighborhood Search, генетический подход с генеративной компонентой и точный поиск для малых экземпляров~\cite{coffman1997,lodi2002,pisinger2010}.

Классические greedy-алгоритмы строят укладку последовательно: объекты сортируются по выбранному приоритету, после чего для каждого перебираются допустимые ориентации и кандидатные позиции. В проекте такая логика реализована через extreme points и gravity drop. Метод быстр и удобен как baseline, но плохо исправляет ранние неудачные решения: если крупный или хрупкий товар размещен неудачно, последующие объекты уже подстраиваются под эту структуру.

Послойный MaxRects-подход использует идею двумерной упаковки прямоугольников внутри слоя. В проекте он адаптирован для паллетизации: выбирается высота слоя, основания коробов размещаются на плоскости, затем алгоритм переходит к следующему уровню. Такой решатель быстро достигает высокой утилизации объема, однако из-за послойной структуры хуже учитывает взаимное расположение хрупких и тяжелых товаров.

Large Neighborhood Search устраняет главный недостаток greedy-схем: вместо локального продолжения уже построенной укладки он периодически удаляет часть размещенных коробов и восстанавливает решение заново. В проекте LNS использует greedy-фазу для старта, несколько destroy-операторов, repair на основе beam search и принятие кандидатов с элементами simulated annealing. По сохраненным результатам именно LNS дает лучший средний итоговый score на общих наборах.

Генетические и ML-ориентированные методы рассматриваются как способ расширить пространство поиска. Хромосома задает порядок объектов и выбор ориентаций, а декодер строит конкретную укладку. Реализованный вариант GAN+GA не использует тяжелые фреймворки обучения: генеративная часть написана на NumPy и влияет на разнообразие популяции и fitness. Такой подход интересен как направление развития, но в текущих экспериментах уступает LNS по среднему score~\cite{zhang2024}.

Точный поиск и branch-and-bound используются только для малых задач из каталога \code{brute\_force\_vs\_algo}. Они полезны как исследовательская точка сравнения: позволяют проверить, насколько эвристические решения близки к более полному перебору. Для рабочих размеров такой подход становится слишком дорогим по времени, поэтому основной практический акцент сделан на LNS и быстрых базовых решателях.
